\chapter*{TÓM TẮT}
\label{Abstract}
\addcontentsline{toc}{chapter}{TÓM TẮT}

Nhu cầu tính toán ngày càng tăng của các Mạng Nơ-ron Sâu (DNN) đã bộc lộ những hạn chế của bộ xử lý đa dụng truyền thống, từ đó đặt ra yêu cầu cấp thiết cho các bộ tăng tốc phần cứng chuyên dụng. Tuy nhiên, việc nghiên cứu và phát triển trong lĩnh vực này thường bị hạn chế bởi các công cụ thiết kế độc quyền và IP phụ thuộc vào nhà cung cấp, với chi phí bản quyền cao và khả năng tùy biến hạn chế. Sự ra đời của kiến trúc tập lệnh (ISA) mã nguồn mở RISC-V đã mang đến một giải pháp thay thế mạnh mẽ và bền vững cho tương lai. Luận văn này trình bày việc thiết kế, triển khai full-stack và đánh giá một Hệ thống trên Vi mạch (SoC) RISC-V có khả năng chạy hệ điều hành Linux, và được tích hợp một bộ tăng tốc DNN chuyên dụng. Mục tiêu chính của luận văn là thực hiện một nghiên cứu toàn diện về hệ sinh thái RISC-V thông qua việc minh chứng một quy trình làm việc mã nguồn mở hoàn chỉnh, từ cấu hình phần cứng mức cao bằng Chisel cho đến việc xây dựng nguyên mẫu trên FPGA. Hệ thống được thiết kế trên nền tảng Chipyard, tích hợp lõi xử lý 64-bit RISC-V Rocket Core với bộ tăng tốc Gemmini có kiến trúc mảng xử lý song song theo nhịp (systolic array) thông qua giao thức kết nối RoCC. Hệ thống SoC đã được triển khai thành công trên bo mạch FPGA Xilinx (AMD) Virtex 7 VC707, khởi chạy thành công hệ điều hành Debian Linux. Kết quả đánh giá hiệu năng cho thấy bộ tăng tốc Gemmini đạt được mức tăng tốc xấp xỉ 1,300 lần trên một tác vụ tính toán đã được tối ưu so với giải pháp phần mềm thuần túy trên Rocket Core. Công trình này đã xác thực thành công nền tảng Chipyard là một nền tảng hoàn thiện để xây dựng các hệ thống RISC-V phức tạp, đồng thời thiết lập một nền tảng full-stack có thể tái tạo cho các nghiên cứu trong tương lai, qua đó cung cấp một lộ trình chuyển đổi từ các quy trình làm việc độc quyền, khép kín sang hệ sinh thái rộng mở cùng tập lệnh RISC-V.

% % Ngữ cảnh, vấn đề
% Trong bối cảnh ngành thiết kế vi mạch, đặc biệt là thiết kế vi mạch số, đang ngày càng phát triển mạnh mẽ, việc hiểu rõ quy trình thiết kế trở nên vô cùng quan trọng.
% % Mục tiêu
% Báo cáo này trình bày quá trình thiết kế và đánh giá một mạch dịch 8-bit theo quy trình chuẩn của thiết kế vi mạch số, bao gồm lập trình Verilog, mô phỏng, tổng hợp mạch, kiểm tra \acrfull{drc} và \acrfull{lvs}.
% % Phương pháp
% Toàn bộ quy trình thiết kế được thực hiện thông qua bộ công cụ của Synopsys bao gồm \acrfull{vcs}, \acrfull{dc}, và \acrfull{icc2}. 
% % Kết quả
% Kết quả cuối cùng là một \acrfull{ip} mạch dịch 8-bit, có thể được sử dụng làm phần tử cơ bản trong các thiết kế lớn hơn.
% % Kết luận
% Bài báo cáo này khẳng định tầm quan trọng của việc nắm vững các quy trình thiết kế và công cụ hiện đại để đáp ứng nhu cầu ngày càng cao trong ngành công nghiệp vi mạch.